% Options for packages loaded elsewhere
\PassOptionsToPackage{unicode}{hyperref}
\PassOptionsToPackage{hyphens}{url}
\documentclass[
]{article}
\usepackage{xcolor}
\usepackage{amsmath,amssymb}
\setcounter{secnumdepth}{-\maxdimen} % remove section numbering
\usepackage{iftex}
\ifPDFTeX
  \usepackage[T1]{fontenc}
  \usepackage[utf8]{inputenc}
  \usepackage{textcomp} % provide euro and other symbols
\else % if luatex or xetex
  \usepackage{unicode-math} % this also loads fontspec
  \defaultfontfeatures{Scale=MatchLowercase}
  \defaultfontfeatures[\rmfamily]{Ligatures=TeX,Scale=1}
\fi
\usepackage{lmodern}
\ifPDFTeX\else
  % xetex/luatex font selection
\fi
% Use upquote if available, for straight quotes in verbatim environments
\IfFileExists{upquote.sty}{\usepackage{upquote}}{}
\IfFileExists{microtype.sty}{% use microtype if available
  \usepackage[]{microtype}
  \UseMicrotypeSet[protrusion]{basicmath} % disable protrusion for tt fonts
}{}
\makeatletter
\@ifundefined{KOMAClassName}{% if non-KOMA class
  \IfFileExists{parskip.sty}{%
    \usepackage{parskip}
  }{% else
    \setlength{\parindent}{0pt}
    \setlength{\parskip}{6pt plus 2pt minus 1pt}}
}{% if KOMA class
  \KOMAoptions{parskip=half}}
\makeatother
\usepackage{color}
\usepackage{fancyvrb}
\newcommand{\VerbBar}{|}
\newcommand{\VERB}{\Verb[commandchars=\\\{\}]}
\DefineVerbatimEnvironment{Highlighting}{Verbatim}{commandchars=\\\{\}}
% Add ',fontsize=\small' for more characters per line
\newenvironment{Shaded}{}{}
\newcommand{\AlertTok}[1]{\textcolor[rgb]{1.00,0.00,0.00}{\textbf{#1}}}
\newcommand{\AnnotationTok}[1]{\textcolor[rgb]{0.38,0.63,0.69}{\textbf{\textit{#1}}}}
\newcommand{\AttributeTok}[1]{\textcolor[rgb]{0.49,0.56,0.16}{#1}}
\newcommand{\BaseNTok}[1]{\textcolor[rgb]{0.25,0.63,0.44}{#1}}
\newcommand{\BuiltInTok}[1]{\textcolor[rgb]{0.00,0.50,0.00}{#1}}
\newcommand{\CharTok}[1]{\textcolor[rgb]{0.25,0.44,0.63}{#1}}
\newcommand{\CommentTok}[1]{\textcolor[rgb]{0.38,0.63,0.69}{\textit{#1}}}
\newcommand{\CommentVarTok}[1]{\textcolor[rgb]{0.38,0.63,0.69}{\textbf{\textit{#1}}}}
\newcommand{\ConstantTok}[1]{\textcolor[rgb]{0.53,0.00,0.00}{#1}}
\newcommand{\ControlFlowTok}[1]{\textcolor[rgb]{0.00,0.44,0.13}{\textbf{#1}}}
\newcommand{\DataTypeTok}[1]{\textcolor[rgb]{0.56,0.13,0.00}{#1}}
\newcommand{\DecValTok}[1]{\textcolor[rgb]{0.25,0.63,0.44}{#1}}
\newcommand{\DocumentationTok}[1]{\textcolor[rgb]{0.73,0.13,0.13}{\textit{#1}}}
\newcommand{\ErrorTok}[1]{\textcolor[rgb]{1.00,0.00,0.00}{\textbf{#1}}}
\newcommand{\ExtensionTok}[1]{#1}
\newcommand{\FloatTok}[1]{\textcolor[rgb]{0.25,0.63,0.44}{#1}}
\newcommand{\FunctionTok}[1]{\textcolor[rgb]{0.02,0.16,0.49}{#1}}
\newcommand{\ImportTok}[1]{\textcolor[rgb]{0.00,0.50,0.00}{\textbf{#1}}}
\newcommand{\InformationTok}[1]{\textcolor[rgb]{0.38,0.63,0.69}{\textbf{\textit{#1}}}}
\newcommand{\KeywordTok}[1]{\textcolor[rgb]{0.00,0.44,0.13}{\textbf{#1}}}
\newcommand{\NormalTok}[1]{#1}
\newcommand{\OperatorTok}[1]{\textcolor[rgb]{0.40,0.40,0.40}{#1}}
\newcommand{\OtherTok}[1]{\textcolor[rgb]{0.00,0.44,0.13}{#1}}
\newcommand{\PreprocessorTok}[1]{\textcolor[rgb]{0.74,0.48,0.00}{#1}}
\newcommand{\RegionMarkerTok}[1]{#1}
\newcommand{\SpecialCharTok}[1]{\textcolor[rgb]{0.25,0.44,0.63}{#1}}
\newcommand{\SpecialStringTok}[1]{\textcolor[rgb]{0.73,0.40,0.53}{#1}}
\newcommand{\StringTok}[1]{\textcolor[rgb]{0.25,0.44,0.63}{#1}}
\newcommand{\VariableTok}[1]{\textcolor[rgb]{0.10,0.09,0.49}{#1}}
\newcommand{\VerbatimStringTok}[1]{\textcolor[rgb]{0.25,0.44,0.63}{#1}}
\newcommand{\WarningTok}[1]{\textcolor[rgb]{0.38,0.63,0.69}{\textbf{\textit{#1}}}}
\usepackage{longtable,booktabs,array}
\newcounter{none} % for unnumbered tables
\usepackage{calc} % for calculating minipage widths
% Correct order of tables after \paragraph or \subparagraph
\usepackage{etoolbox}
\makeatletter
\patchcmd\longtable{\par}{\if@noskipsec\mbox{}\fi\par}{}{}
\makeatother
% Allow footnotes in longtable head/foot
\IfFileExists{footnotehyper.sty}{\usepackage{footnotehyper}}{\usepackage{footnote}}
\makesavenoteenv{longtable}
\setlength{\emergencystretch}{3em} % prevent overfull lines
\providecommand{\tightlist}{%
  \setlength{\itemsep}{0pt}\setlength{\parskip}{0pt}}
\usepackage{bookmark}
\IfFileExists{xurl.sty}{\usepackage{xurl}}{} % add URL line breaks if available
\urlstyle{same}
\hypersetup{
  hidelinks,
  pdfcreator={LaTeX via pandoc}}

\author{}
\date{}

\begin{document}

\section{Security Audit Report: DID
Contracts}\label{security-audit-report-did-contracts}

\textbf{Audited Contracts:}

\begin{itemize}
\tightlist
\item
  \texttt{contracts/did/Registry.sol}
\item
  \texttt{contracts/did/IdentityNFT.sol}
\end{itemize}

\textbf{Audit Date:} 2026-01-30 \textbf{Auditor:} Trail of Bits-level
Review \textbf{Solidity Version:} \^{}0.8.31

\begin{center}\rule{0.5\linewidth}{0.5pt}\end{center}

\subsection{Executive Summary}\label{executive-summary}

The DID (Decentralized Identifier) contracts implement a multi-chain
identity registry with stake-based registration, NFT-bound ownership,
and delegation capabilities. The audit identified \textbf{2 Critical},
\textbf{4 High}, \textbf{5 Medium}, \textbf{4 Low}, and \textbf{6
Informational} issues.

\subsubsection{Findings Summary}\label{findings-summary}

{\def\LTcaptype{none} % do not increment counter
\begin{longtable}[]{@{}lll@{}}
\toprule\noalign{}
Severity & Count & Status \\
\midrule\noalign{}
\endhead
\bottomrule\noalign{}
\endlastfoot
Critical & 2 & Must Fix \\
High & 4 & Must Fix \\
Medium & 5 & Should Fix \\
Low & 4 & Consider Fixing \\
Informational & 6 & Best Practice \\
\end{longtable}
}

\begin{center}\rule{0.5\linewidth}{0.5pt}\end{center}

\subsection{Critical Issues}\label{critical-issues}

\subsubsection{\texorpdfstring{C-01: Reentrancy in \texttt{unclaim()} -
Stake Returned Before State
Cleared}{C-01: Reentrancy in unclaim() - Stake Returned Before State Cleared}}\label{c-01-reentrancy-in-unclaim---stake-returned-before-state-cleared}

\textbf{File:} \texttt{Registry.sol:251-272}

\textbf{Description:} The \texttt{unclaim()} function transfers staked
tokens to the caller before clearing state. If the staking token is a
malicious ERC20 with a callback (e.g., ERC777), an attacker can re-enter
and exploit inconsistent state.

\begin{Shaded}
\begin{Highlighting}[]
\NormalTok{function unclaim(string calldata did) external \{}
\NormalTok{    \_requireOwner(did);}

\NormalTok{    IdentityData storage data = \_data[did];}
\NormalTok{    uint256 nftId = data.boundNft;}

\NormalTok{    // VULNERABILITY: External call before state changes}
\NormalTok{    if (data.stakedTokens \textgreater{} 0) \{}
\NormalTok{        stakingToken.transfer(msg.sender, data.stakedTokens);  // \textless{}{-}{-} External call}
\NormalTok{    \}}

\NormalTok{    identityNft.burn(nftId);  // \textless{}{-}{-} Another external call}

\NormalTok{    // State cleared AFTER external calls}
\NormalTok{    delete \_owners[did];}
\NormalTok{    delete tokenToDID[nftId];}
\NormalTok{    delete \_data[did];}
\NormalTok{    delete \_delegatees[did];}
\NormalTok{\}}
\end{Highlighting}
\end{Shaded}

\textbf{Attack Scenario:}

\begin{enumerate}
\def\labelenumi{\arabic{enumi}.}
\tightlist
\item
  Attacker claims DID with malicious ERC777-like token as staking token
\item
  Attacker calls \texttt{unclaim()}
\item
  On \texttt{transfer()}, token triggers callback to attacker
\item
  Attacker re-enters via another function while
  \texttt{\_owners{[}did{]}} still points to them
\item
  Attacker can manipulate data or double-spend
\end{enumerate}

\textbf{Recommendation:} Apply Checks-Effects-Interactions pattern:

\begin{Shaded}
\begin{Highlighting}[]
\NormalTok{function unclaim(string calldata did) external \{}
\NormalTok{    \_requireOwner(did);}

\NormalTok{    IdentityData storage data = \_data[did];}
\NormalTok{    uint256 nftId = data.boundNft;}
\NormalTok{    uint256 stakeToReturn = data.stakedTokens;}

\NormalTok{    // Effects FIRST}
\NormalTok{    delete \_owners[did];}
\NormalTok{    delete tokenToDID[nftId];}
\NormalTok{    delete \_data[did];}
\NormalTok{    delete \_delegatees[did];}

\NormalTok{    emit IdentityUnclaimed(did, nftId);}

\NormalTok{    // Interactions LAST}
\NormalTok{    identityNft.burn(nftId);}
\NormalTok{    if (stakeToReturn \textgreater{} 0) \{}
\NormalTok{        stakingToken.transfer(msg.sender, stakeToReturn);}
\NormalTok{    \}}
\NormalTok{\}}
\end{Highlighting}
\end{Shaded}

Also consider adding \texttt{ReentrancyGuard} from OpenZeppelin.

\begin{center}\rule{0.5\linewidth}{0.5pt}\end{center}

\subsubsection{C-02: NFT Transfer Desynchronizes DID
Ownership}\label{c-02-nft-transfer-desynchronizes-did-ownership}

\textbf{File:} \texttt{Registry.sol:229-230},
\texttt{IdentityNFT.sol:64-69}

\textbf{Description:} The \texttt{\_owners{[}did{]}} mapping in Registry
and actual NFT ownership can become desynchronized. The NFT is
transferable, but when transferred, the Registry\textquotesingle s
\texttt{\_owners} mapping is not updated.

\begin{Shaded}
\begin{Highlighting}[]
\NormalTok{// Registry stores owner separately}
\NormalTok{\_owners[did] = params.owner;}

\NormalTok{// IdentityNFT is a standard ERC721 {-} fully transferable}
\NormalTok{\_safeMint(to, tokenId);}
\end{Highlighting}
\end{Shaded}

\textbf{Attack Scenario:}

\begin{enumerate}
\def\labelenumi{\arabic{enumi}.}
\tightlist
\item
  Alice claims \texttt{did:lux:alice}, becomes owner in both Registry
  and NFT
\item
  Alice transfers NFT to Bob via standard ERC721 \texttt{transferFrom()}
\item
  Registry still shows \texttt{\_owners{[}"did:lux:alice"{]}\ =\ Alice}
\item
  Alice can still call \texttt{unclaim()}, \texttt{setKeys()}, etc. even
  though Bob owns the NFT
\item
  Alice steals staked tokens; Bob loses NFT when Alice calls
  \texttt{unclaim()}
\end{enumerate}

\textbf{Impact:} Complete identity theft, stake drainage, unauthorized
identity modifications.

\textbf{Recommendation:}

Option A - Override NFT transfer to sync Registry:

\begin{Shaded}
\begin{Highlighting}[]
\NormalTok{// In IdentityNFT.sol}
\NormalTok{function \_update(address to, uint256 tokenId, address auth) internal override returns (address) \{}
\NormalTok{    address from = super.\_update(to, tokenId, auth);}
\NormalTok{    if (from != address(0) \&\& to != address(0)) \{}
\NormalTok{        // Notify registry of ownership change}
\NormalTok{        IRegistry(registry).syncOwnership(tokenId, to);}
\NormalTok{    \}}
\NormalTok{    return from;}
\NormalTok{\}}
\end{Highlighting}
\end{Shaded}

Option B - Make \texttt{\_requireOwner()} check NFT ownership:

\begin{Shaded}
\begin{Highlighting}[]
\NormalTok{function \_requireOwner(string memory did) internal view \{}
\NormalTok{    IdentityData storage data = \_data[did];}
\NormalTok{    address nftOwner = identityNft.ownerOf(data.boundNft);}
\NormalTok{    if (nftOwner != msg.sender) revert Unauthorized();}
\NormalTok{\}}
\end{Highlighting}
\end{Shaded}

Option B is simpler and recommended. The \texttt{\_owners} mapping can
be deprecated in favor of NFT ownership as the source of truth.

\begin{center}\rule{0.5\linewidth}{0.5pt}\end{center}

\subsection{High Severity Issues}\label{high-severity-issues}

\subsubsection{H-01: Missing Return Value Check on ERC20
Transfers}\label{h-01-missing-return-value-check-on-erc20-transfers}

\textbf{File:} \texttt{Registry.sol:224,\ 259,\ 320}

\textbf{Description:} The contract calls \texttt{transferFrom()} and
\texttt{transfer()} without checking return values. Some ERC20 tokens
(USDT, BNB) return \texttt{false} instead of reverting on failure.

\begin{Shaded}
\begin{Highlighting}[]
\NormalTok{// Line 224 {-} No return check}
\NormalTok{stakingToken.transferFrom(msg.sender, address(this), params.stakeAmount);}

\NormalTok{// Line 259 {-} No return check}
\NormalTok{stakingToken.transfer(msg.sender, data.stakedTokens);}

\NormalTok{// Line 320 {-} No return check}
\NormalTok{stakingToken.transferFrom(msg.sender, address(this), amount);}
\end{Highlighting}
\end{Shaded}

\textbf{Impact:} Silent transfer failures leading to:

\begin{itemize}
\tightlist
\item
  Claims without actual stake being deposited
\item
  State inconsistency between recorded stake and actual balance
\end{itemize}

\textbf{Recommendation:} Use OpenZeppelin\textquotesingle s
\texttt{SafeERC20}:

\begin{Shaded}
\begin{Highlighting}[]
\NormalTok{import \{SafeERC20\} from "@openzeppelin/contracts/token/ERC20/utils/SafeERC20.sol";}

\NormalTok{using SafeERC20 for IERC20;}

\NormalTok{// Then use:}
\NormalTok{stakingToken.safeTransferFrom(msg.sender, address(this), params.stakeAmount);}
\NormalTok{stakingToken.safeTransfer(msg.sender, data.stakedTokens);}
\end{Highlighting}
\end{Shaded}

\begin{center}\rule{0.5\linewidth}{0.5pt}\end{center}

\subsubsection{H-02: Delegation Data Not Cleared on Unclaim - Delegation
Accounting
Corruption}\label{h-02-delegation-data-not-cleared-on-unclaim---delegation-accounting-corruption}

\textbf{File:} \texttt{Registry.sol:269}

\textbf{Description:} The \texttt{unclaim()} function deletes
\texttt{\_delegatees{[}did{]}} but does NOT clear the individual
\texttt{delegations{[}did{]}{[}delegatee{]}} mappings. This leaves
orphaned delegation records.

\begin{Shaded}
\begin{Highlighting}[]
\NormalTok{// Line 269 {-} Only clears the array, not the mapping entries}
\NormalTok{delete \_delegatees[did];}

\NormalTok{// delegations[did][*] entries remain!}
\end{Highlighting}
\end{Shaded}

\textbf{Impact:}

\begin{itemize}
\tightlist
\item
  If the same DID is later reclaimed, old delegation amounts may be
  accessible
\item
  Corrupts delegation accounting for governance calculations
\item
  Could be exploited to resurrect old voting power
\end{itemize}

\textbf{Recommendation:} Clear delegation mappings before deleting the
array:

\begin{Shaded}
\begin{Highlighting}[]
\NormalTok{string[] memory dels = \_delegatees[did];}
\NormalTok{for (uint256 i = 0; i \textless{} dels.length; i++) \{}
\NormalTok{    delete delegations[did][dels[i]];}
\NormalTok{\}}
\NormalTok{delete \_delegatees[did];}
\end{Highlighting}
\end{Shaded}

\begin{center}\rule{0.5\linewidth}{0.5pt}\end{center}

\subsubsection{\texorpdfstring{H-03: Front-Running on \texttt{claim()} -
Name
Squatting}{H-03: Front-Running on claim() - Name Squatting}}\label{h-03-front-running-on-claim---name-squatting}

\textbf{File:} \texttt{Registry.sol:205-245}

\textbf{Description:} The \texttt{claim()} function is vulnerable to
front-running. An attacker monitoring the mempool can see a pending
claim transaction and submit their own with higher gas price to steal
the desired name.

\textbf{Attack Scenario:}

\begin{enumerate}
\def\labelenumi{\arabic{enumi}.}
\tightlist
\item
  Alice submits \texttt{claim("valuable\_name")}
\item
  Attacker sees transaction in mempool
\item
  Attacker submits same claim with higher gas price
\item
  Attacker\textquotesingle s transaction executes first, stealing the
  name
\item
  Alice\textquotesingle s transaction reverts with
  \texttt{IdentityNotAvailable}
\end{enumerate}

\textbf{Impact:} Valuable names can be systematically front-run and held
for ransom.

\textbf{Recommendation:} Implement commit-reveal scheme:

\begin{Shaded}
\begin{Highlighting}[]
\NormalTok{mapping(bytes32 =\textgreater{} uint256) public commitments;}

\NormalTok{function commit(bytes32 commitment) external \{}
\NormalTok{    commitments[commitment] = block.timestamp;}
\NormalTok{\}}

\NormalTok{function claim(ClaimParams calldata params, bytes32 secret) external \{}
\NormalTok{    bytes32 commitment = keccak256(abi.encodePacked(params.name, params.owner, secret));}
\NormalTok{    require(commitments[commitment] != 0, "No commitment");}
\NormalTok{    require(block.timestamp \textgreater{}= commitments[commitment] + 1 minutes, "Too early");}
\NormalTok{    require(block.timestamp \textless{}= commitments[commitment] + 24 hours, "Expired");}

\NormalTok{    delete commitments[commitment];}
\NormalTok{    // ... rest of claim logic}
\NormalTok{\}}
\end{Highlighting}
\end{Shaded}

\begin{center}\rule{0.5\linewidth}{0.5pt}\end{center}

\subsubsection{\texorpdfstring{H-04: \texttt{\_validName()} Allows
Confusable Characters - Homograph
Attacks}{H-04: \_validName() Allows Confusable Characters - Homograph Attacks}}\label{h-04-_validname-allows-confusable-characters---homograph-attacks}

\textbf{File:} \texttt{Registry.sol:435-448}

\textbf{Description:} The name validation allows uppercase letters,
enabling homograph attacks with visually similar names:

\begin{Shaded}
\begin{Highlighting}[]
\NormalTok{bool valid = (c \textgreater{}= 0x30 \&\& c \textless{}= 0x39) ||  // 0{-}9}
\NormalTok{            (c \textgreater{}= 0x41 \&\& c \textless{}= 0x5A) ||   // A{-}Z  \textless{}{-}{-} Uppercase allowed}
\NormalTok{            (c \textgreater{}= 0x61 \&\& c \textless{}= 0x7A) ||   // a{-}z}
\NormalTok{            (c == 0x5F);                   // \_}
\end{Highlighting}
\end{Shaded}

\textbf{Attack Scenario:}

\begin{itemize}
\tightlist
\item
  Attacker registers \texttt{ALICE} (all uppercase)
\item
  User thinks they\textquotesingle re interacting with \texttt{alice}
\item
  Two different DIDs: \texttt{did:lux:ALICE} vs \texttt{did:lux:alice}
\item
  Phishing and confusion attacks possible
\end{itemize}

Additional concerns:

\begin{itemize}
\tightlist
\item
  \texttt{\_} can be confused with \texttt{-} or space
\item
  \texttt{l} (lowercase L) vs \texttt{1} (one) vs \texttt{I} (uppercase
  i)
\item
  \texttt{0} (zero) vs \texttt{O} (uppercase o)
\end{itemize}

\textbf{Recommendation:} Enforce lowercase-only names and normalize on
storage:

\begin{Shaded}
\begin{Highlighting}[]
\NormalTok{function \_validName(string calldata name) internal pure returns (bool) \{}
\NormalTok{    bytes memory b = bytes(name);}
\NormalTok{    if (b.length == 0 || b.length \textgreater{} 63) return false;}

\NormalTok{    for (uint256 i = 0; i \textless{} b.length; i++) \{}
\NormalTok{        bytes1 c = b[i];}
\NormalTok{        // Only lowercase and digits}
\NormalTok{        bool valid = (c \textgreater{}= 0x30 \&\& c \textless{}= 0x39) ||  // 0{-}9}
\NormalTok{                    (c \textgreater{}= 0x61 \&\& c \textless{}= 0x7A);      // a{-}z only}
\NormalTok{        if (!valid) return false;}
\NormalTok{    \}}
\NormalTok{    return true;}
\NormalTok{\}}
\end{Highlighting}
\end{Shaded}

\begin{center}\rule{0.5\linewidth}{0.5pt}\end{center}

\subsection{Medium Severity Issues}\label{medium-severity-issues}

\subsubsection{\texorpdfstring{M-01: \texttt{IdentityNFT.mint()} Uses
\texttt{\_safeMint()} - Callback
Risk}{M-01: IdentityNFT.mint() Uses \_safeMint() - Callback Risk}}\label{m-01-identitynftmint-uses-_safemint---callback-risk}

\textbf{File:} \texttt{IdentityNFT.sol:68}

\textbf{Description:} The \texttt{mint()} function uses
\texttt{\_safeMint()} which calls \texttt{onERC721Received()} on the
recipient if it\textquotesingle s a contract. This external call happens
before the Registry completes its state updates in \texttt{claim()}.

\begin{Shaded}
\begin{Highlighting}[]
\NormalTok{// In IdentityNFT.mint()}
\NormalTok{\_safeMint(to, tokenId);  // Calls external contract}

\NormalTok{// Back in Registry.claim(), state updates happen AFTER mint returns}
\NormalTok{\_owners[did] = params.owner;}
\NormalTok{tokenToDID[nftId] = did;}
\NormalTok{\_data[did] = IdentityData(\{...\});}
\end{Highlighting}
\end{Shaded}

\textbf{Impact:} A malicious recipient contract could re-enter the
Registry during the callback with inconsistent state.

\textbf{Recommendation:} Either use \texttt{\_mint()} instead of
\texttt{\_safeMint()}, or ensure Registry state is set before calling
\texttt{identityNft.mint()}.

\begin{center}\rule{0.5\linewidth}{0.5pt}\end{center}

\subsubsection{M-02: No Stake Decrease Function - Locked Tokens Beyond
Minimum}\label{m-02-no-stake-decrease-function---locked-tokens-beyond-minimum}

\textbf{File:} \texttt{Registry.sol:317-327}

\textbf{Description:} Users can only increase stake via
\texttt{increaseStake()}. There is no mechanism to withdraw excess stake
above the minimum requirement.

\begin{Shaded}
\begin{Highlighting}[]
\NormalTok{function increaseStake(string calldata did, uint256 amount) external \{}
\NormalTok{    \_requireOwner(did);}
\NormalTok{    stakingToken.transferFrom(msg.sender, address(this), amount);}
\NormalTok{    IdentityData storage data = \_data[did];}
\NormalTok{    data.stakedTokens += amount;}
\NormalTok{    // ... no way to decrease}
\NormalTok{\}}
\end{Highlighting}
\end{Shaded}

\textbf{Impact:} Users who overstake (intentionally or accidentally) can
only recover tokens by completely unclaiming their identity.

\textbf{Recommendation:} Add \texttt{decreaseStake()} function with
minimum stake enforcement:

\begin{Shaded}
\begin{Highlighting}[]
\NormalTok{function decreaseStake(string calldata did, uint256 amount) external \{}
\NormalTok{    \_requireOwner(did);}
\NormalTok{    IdentityData storage data = \_data[did];}

\NormalTok{    // Extract name from DID for requirement calculation}
\NormalTok{    uint256 minimum = stakeRequirement(...);}
\NormalTok{    require(data.stakedTokens {-} amount \textgreater{}= minimum, "Below minimum");}

\NormalTok{    data.stakedTokens {-}= amount;}
\NormalTok{    stakingToken.transfer(msg.sender, amount);}
\NormalTok{\}}
\end{Highlighting}
\end{Shaded}

\begin{center}\rule{0.5\linewidth}{0.5pt}\end{center}

\subsubsection{\texorpdfstring{M-03: \texttt{\_delegatees} Array Can
Grow Unbounded - DoS
Vector}{M-03: \_delegatees Array Can Grow Unbounded - DoS Vector}}\label{m-03-_delegatees-array-can-grow-unbounded---dos-vector}

\textbf{File:} \texttt{Registry.sol:101}

\textbf{Description:} The \texttt{\_delegatees} array has no size limit.
While there\textquotesingle s no \texttt{delegate()} function visible in
the audited code, the struct and storage exist. If delegation is
implemented, unbounded array growth could cause gas limit issues.

\textbf{Impact:} Potential denial-of-service if delegation is added
without bounds.

\textbf{Recommendation:}

\begin{itemize}
\tightlist
\item
  Implement maximum delegatee limit (e.g., 100)
\item
  Use enumerable set pattern for O(1) removal
\item
  Consider removing unused delegation storage if not needed
\end{itemize}

\begin{center}\rule{0.5\linewidth}{0.5pt}\end{center}

\subsubsection{M-04: Owner Can Change Pricing to Zero - Free Identity
Claims}\label{m-04-owner-can-change-pricing-to-zero---free-identity-claims}

\textbf{File:} \texttt{Registry.sol:410-422}

\textbf{Description:} The \texttt{setPricing()} function allows setting
prices to zero with no lower bound validation:

\begin{Shaded}
\begin{Highlighting}[]
\NormalTok{function setPricing(}
\NormalTok{    uint256 p1, uint256 p2, uint256 p3, uint256 p4, uint256 p5,}
\NormalTok{    uint256 discountBps}
\NormalTok{) external onlyOwner \{}
\NormalTok{    require(discountBps \textless{}= 10000, "Invalid discount");}
\NormalTok{    // No minimum price validation}
\NormalTok{    price1Char = p1;  // Can be 0}
\NormalTok{    // ...}
\NormalTok{\}}
\end{Highlighting}
\end{Shaded}

\textbf{Impact:} Admin error or compromised owner could set prices to 0,
enabling mass squatting of valuable names.

\textbf{Recommendation:} Add minimum price validation:

\begin{Shaded}
\begin{Highlighting}[]
\NormalTok{require(p1 \textgreater{} 0 \&\& p2 \textgreater{} 0 \&\& p3 \textgreater{} 0 \&\& p4 \textgreater{} 0 \&\& p5 \textgreater{} 0, "Invalid price");}
\end{Highlighting}
\end{Shaded}

\begin{center}\rule{0.5\linewidth}{0.5pt}\end{center}

\subsubsection{\texorpdfstring{M-05: \texttt{records} Mapping Cleanup
Missing on
Unclaim}{M-05: records Mapping Cleanup Missing on Unclaim}}\label{m-05-records-mapping-cleanup-missing-on-unclaim}

\textbf{File:} \texttt{Registry.sol:95,\ 266-268}

\textbf{Description:} The \texttt{records{[}did{]}} mapping is not
cleared when a DID is unclaimed:

\begin{Shaded}
\begin{Highlighting}[]
\NormalTok{// Line 95}
\NormalTok{mapping(string =\textgreater{} mapping(string =\textgreater{} string)) public records;}

\NormalTok{// Line 266{-}268 {-} records NOT deleted}
\NormalTok{delete \_owners[did];}
\NormalTok{delete tokenToDID[nftId];}
\NormalTok{delete \_data[did];}
\NormalTok{delete \_delegatees[did];}
\NormalTok{// Missing: records[did] not cleared}
\end{Highlighting}
\end{Shaded}

\textbf{Impact:}

\begin{itemize}
\tightlist
\item
  Old records persist and may be associated with a new owner if DID is
  reclaimed
\item
  Privacy leak - old data remains accessible
\item
  Storage bloat
\end{itemize}

\textbf{Recommendation:} Note: Clearing a nested mapping is expensive.
Consider:

\begin{enumerate}
\def\labelenumi{\arabic{enumi}.}
\tightlist
\item
  Version number per DID to invalidate old records
\item
  Explicit key tracking for cleanup
\item
  Accept the limitation and document it
\end{enumerate}

\begin{center}\rule{0.5\linewidth}{0.5pt}\end{center}

\subsection{Low Severity Issues}\label{low-severity-issues}

\subsubsection{\texorpdfstring{L-01: Missing Events in
\texttt{IdentityNFT}}{L-01: Missing Events in IdentityNFT}}\label{l-01-missing-events-in-identitynft}

\textbf{File:} \texttt{IdentityNFT.sol:64-78}

\textbf{Description:} The \texttt{mint()} and \texttt{burn()} functions
don\textquotesingle t emit custom events. While ERC721 emits
\texttt{Transfer} events, custom events would improve off-chain
indexing.

\textbf{Recommendation:} Add dedicated events:

\begin{Shaded}
\begin{Highlighting}[]
\NormalTok{event IdentityMinted(address indexed to, uint256 indexed tokenId);}
\NormalTok{event IdentityBurned(uint256 indexed tokenId);}
\end{Highlighting}
\end{Shaded}

\begin{center}\rule{0.5\linewidth}{0.5pt}\end{center}

\subsubsection{\texorpdfstring{L-02: \texttt{IdentityNFT} Registry Can
Be Set to Zero After Initial
Setup}{L-02: IdentityNFT Registry Can Be Set to Zero After Initial Setup}}\label{l-02-identitynft-registry-can-be-set-to-zero-after-initial-setup}

\textbf{File:} \texttt{IdentityNFT.sol:46-49}

\textbf{Description:} The \texttt{setRegistry()} function validates
against zero address, but can be called multiple times. Once set,
changing the registry could break the system.

\begin{Shaded}
\begin{Highlighting}[]
\NormalTok{function setRegistry(address registry\_) external onlyOwner \{}
\NormalTok{    if (registry\_ == address(0)) revert ZeroAddress();}
\NormalTok{    registry = registry\_;  // Can be changed after deployment}
\NormalTok{\}}
\end{Highlighting}
\end{Shaded}

\textbf{Recommendation:} Make registry immutable after first set, or
remove setter entirely:

\begin{Shaded}
\begin{Highlighting}[]
\NormalTok{function setRegistry(address registry\_) external onlyOwner \{}
\NormalTok{    if (registry != address(0)) revert RegistryAlreadySet();}
\NormalTok{    if (registry\_ == address(0)) revert ZeroAddress();}
\NormalTok{    registry = registry\_;}
\NormalTok{\}}
\end{Highlighting}
\end{Shaded}

\begin{center}\rule{0.5\linewidth}{0.5pt}\end{center}

\subsubsection{L-03: No Timelock on Admin
Functions}\label{l-03-no-timelock-on-admin-functions}

\textbf{File:} \texttt{Registry.sol:402-422}

\textbf{Description:} Critical admin functions like \texttt{setChain()}
and \texttt{setPricing()} execute immediately without timelock. Users
have no warning before parameter changes.

\textbf{Recommendation:} Implement timelock for sensitive operations or
use OpenZeppelin\textquotesingle s TimelockController.

\begin{center}\rule{0.5\linewidth}{0.5pt}\end{center}

\subsubsection{\texorpdfstring{L-04: \texttt{claim()} Allows
Zero-Address
Owner}{L-04: claim() Allows Zero-Address Owner}}\label{l-04-claim-allows-zero-address-owner}

\textbf{File:} \texttt{Registry.sol:227}

\textbf{Description:} The \texttt{params.owner} is not validated before
minting NFT:

\begin{Shaded}
\begin{Highlighting}[]
\NormalTok{uint256 nftId = identityNft.mint(params.owner);  // params.owner could be address(0)}
\end{Highlighting}
\end{Shaded}

While \texttt{IdentityNFT.mint()} checks for zero address, the revert
message would be confusing.

\textbf{Recommendation:} Validate owner in Registry before calling mint:

\begin{Shaded}
\begin{Highlighting}[]
\NormalTok{if (params.owner == address(0)) revert ZeroAddress();}
\end{Highlighting}
\end{Shaded}

\begin{center}\rule{0.5\linewidth}{0.5pt}\end{center}

\subsection{Informational / Gas
Optimizations}\label{informational--gas-optimizations}

\subsubsection{I-01: String Comparison Gas
Optimization}\label{i-01-string-comparison-gas-optimization}

\textbf{File:} \texttt{Registry.sol:210,\ 216,\ 219}

\textbf{Description:} Multiple string operations use \texttt{bytes()}
conversion. Consider caching:

\begin{Shaded}
\begin{Highlighting}[]
\NormalTok{// Current {-} multiple conversions}
\NormalTok{if (bytes(methods[params.chainId]).length == 0) revert InvalidChain(params.chainId);}
\end{Highlighting}
\end{Shaded}

\textbf{Recommendation:} Cache method once and reuse.

\begin{center}\rule{0.5\linewidth}{0.5pt}\end{center}

\subsubsection{I-02: Use Custom Errors
Throughout}\label{i-02-use-custom-errors-throughout}

\textbf{File:} \texttt{Registry.sol:414}

\textbf{Description:} Mix of \texttt{require} with string and custom
errors:

\begin{Shaded}
\begin{Highlighting}[]
\NormalTok{require(discountBps \textless{}= 10000, "Invalid discount");  // String error}
\end{Highlighting}
\end{Shaded}

\textbf{Recommendation:} Use custom error for consistency:

\begin{Shaded}
\begin{Highlighting}[]
\NormalTok{error InvalidDiscount();}
\NormalTok{if (discountBps \textgreater{} 10000) revert InvalidDiscount();}
\end{Highlighting}
\end{Shaded}

\begin{center}\rule{0.5\linewidth}{0.5pt}\end{center}

\subsubsection{\texorpdfstring{I-03: Consider \texttt{bytes32} for
DIDs}{I-03: Consider bytes32 for DIDs}}\label{i-03-consider-bytes32-for-dids}

\textbf{Description:} Storing DIDs as strings is expensive. For
fixed-format DIDs, consider hashing to \texttt{bytes32}:

\begin{Shaded}
\begin{Highlighting}[]
\NormalTok{mapping(bytes32 =\textgreater{} address) private \_owners;  // Hash of DID}

\NormalTok{function \_didHash(string memory did) internal pure returns (bytes32) \{}
\NormalTok{    return keccak256(bytes(did));}
\NormalTok{\}}
\end{Highlighting}
\end{Shaded}

\begin{center}\rule{0.5\linewidth}{0.5pt}\end{center}

\subsubsection{\texorpdfstring{I-04: Unused \texttt{Delegation}
Struct}{I-04: Unused Delegation Struct}}\label{i-04-unused-delegation-struct}

\textbf{File:} \texttt{Registry.sol:64-67}

\textbf{Description:} The \texttt{Delegation} struct is defined but only
used in event parameter. The actual delegation storage uses separate
mappings.

\textbf{Recommendation:} Remove if unused, or restructure delegation
storage to use the struct.

\begin{center}\rule{0.5\linewidth}{0.5pt}\end{center}

\subsubsection{I-05: Missing NatSpec
Documentation}\label{i-05-missing-natspec-documentation}

\textbf{File:} Both contracts

\textbf{Description:} Several functions lack NatSpec documentation,
particularly:

\begin{itemize}
\tightlist
\item
  \texttt{\_authorizeUpgrade()}
\item
  Internal functions
\item
  Event parameters
\end{itemize}

\textbf{Recommendation:} Add comprehensive NatSpec for all
public/external functions.

\begin{center}\rule{0.5\linewidth}{0.5pt}\end{center}

\subsubsection{I-06: Token Counter Could Use Unchecked
Increment}\label{i-06-token-counter-could-use-unchecked-increment}

\textbf{File:} \texttt{IdentityNFT.sol:67}

\textbf{Description:} The token counter increment can use unchecked math
since overflow is practically impossible:

\begin{Shaded}
\begin{Highlighting}[]
\NormalTok{// Current}
\NormalTok{uint256 tokenId = \_tokenIdCounter++;}

\NormalTok{// Optimized}
\NormalTok{uint256 tokenId;}
\NormalTok{unchecked \{ tokenId = \_tokenIdCounter++; \}}
\end{Highlighting}
\end{Shaded}

\begin{center}\rule{0.5\linewidth}{0.5pt}\end{center}

\subsection{Upgrade Safety Analysis
(UUPS)}\label{upgrade-safety-analysis-uups}

\subsubsection{Storage Layout}\label{storage-layout}

The Registry contract uses UUPS upgrade pattern correctly:

\begin{itemize}
\tightlist
\item
  \texttt{\_disableInitializers()} in constructor prevents
  implementation initialization
\item
  \texttt{\_authorizeUpgrade()} restricted to owner
\item
  Inherits from OpenZeppelin upgradeable contracts
\end{itemize}

\textbf{Recommendations:}

\begin{enumerate}
\def\labelenumi{\arabic{enumi}.}
\tightlist
\item
  Add storage gap for future upgrades:
\end{enumerate}

\begin{Shaded}
\begin{Highlighting}[]
\NormalTok{uint256[50] private \_\_gap;}
\end{Highlighting}
\end{Shaded}

\begin{enumerate}
\def\labelenumi{\arabic{enumi}.}
\setcounter{enumi}{1}
\tightlist
\item
  Document storage layout for upgrade safety
\item
  Consider adding upgrade timelock
\end{enumerate}

\subsubsection{Storage Collision Risk}\label{storage-collision-risk}

Current storage order is safe, but new variables MUST be added at the
end. Document variable positions:

\begin{Shaded}
\begin{Highlighting}[]
\NormalTok{// Slot 0: stakingToken (inherited)}
\NormalTok{// Slot 1: identityNft}
\NormalTok{// Slot 2: methods mapping}
\NormalTok{// ... etc}
\end{Highlighting}
\end{Shaded}

\begin{center}\rule{0.5\linewidth}{0.5pt}\end{center}

\subsection{Cross-Chain Security
Considerations}\label{cross-chain-security-considerations}

\begin{enumerate}
\def\labelenumi{\arabic{enumi}.}
\item
  \textbf{Chain ID Spoofing:} The contract trusts
  \texttt{params.chainId} input. On L2s or bridges, verify against
  \texttt{block.chainid}.
\item
  \textbf{Replay Protection:} DIDs include chain-specific method
  (did:lux:, did:zoo:), providing natural replay protection.
\item
  \textbf{Oracle Dependencies:} No external oracle dependencies found -
  good.
\item
  \textbf{Bridge Interactions:} If DIDs are bridged, ensure ownership
  proof travels with the message.
\end{enumerate}

\begin{center}\rule{0.5\linewidth}{0.5pt}\end{center}

\subsection{Recommendations Summary}\label{recommendations-summary}

\subsubsection{Must Fix (Critical/High)}\label{must-fix-criticalhigh}

\begin{enumerate}
\def\labelenumi{\arabic{enumi}.}
\tightlist
\item
  \textbf{C-01:} Apply CEI pattern in \texttt{unclaim()}, add
  ReentrancyGuard
\item
  \textbf{C-02:} Derive ownership from NFT, not stored mapping
\item
  \textbf{H-01:} Use SafeERC20 for all token transfers
\item
  \textbf{H-02:} Clear delegation mappings on unclaim
\item
  \textbf{H-03:} Implement commit-reveal for claims
\item
  \textbf{H-04:} Normalize names to lowercase only
\end{enumerate}

\subsubsection{Should Fix (Medium)}\label{should-fix-medium}

\begin{enumerate}
\def\labelenumi{\arabic{enumi}.}
\tightlist
\item
  \textbf{M-01:} Use \texttt{\_mint()} or reorder state updates
\item
  \textbf{M-02:} Add \texttt{decreaseStake()} function
\item
  \textbf{M-03:} Bound delegatee array size
\item
  \textbf{M-04:} Add minimum price validation
\item
  \textbf{M-05:} Handle records cleanup or version
\end{enumerate}

\subsubsection{Consider Fixing (Low)}\label{consider-fixing-low}

\begin{enumerate}
\def\labelenumi{\arabic{enumi}.}
\tightlist
\item
  \textbf{L-01:} Add custom events for mint/burn
\item
  \textbf{L-02:} Make registry immutable after set
\item
  \textbf{L-03:} Add timelock for admin functions
\item
  \textbf{L-04:} Validate owner address in claim
\end{enumerate}

\begin{center}\rule{0.5\linewidth}{0.5pt}\end{center}

\subsection{Test Vectors}\label{test-vectors}

\subsubsection{C-01 Reentrancy Test}\label{c-01-reentrancy-test}

\begin{Shaded}
\begin{Highlighting}[]
\NormalTok{contract MaliciousToken \{}
\NormalTok{    Registry registry;}

\NormalTok{    function transfer(address, uint256) external returns (bool) \{}
\NormalTok{        // Re{-}enter while owner still set}
\NormalTok{        registry.setKeys("did:lux:attacker", "evil", "evil");}
\NormalTok{        return true;}
\NormalTok{    \}}
\NormalTok{\}}
\end{Highlighting}
\end{Shaded}

\subsubsection{C-02 Ownership Desync
Test}\label{c-02-ownership-desync-test}

\begin{Shaded}
\begin{Highlighting}[]
\NormalTok{function testOwnershipDesync() public \{}
\NormalTok{    // Alice claims}
\NormalTok{    registry.claim(params);}

\NormalTok{    // Alice transfers NFT to Bob}
\NormalTok{    identityNft.transferFrom(alice, bob, tokenId);}

\NormalTok{    // Alice can still unclaim and steal stake}
\NormalTok{    vm.prank(alice);}
\NormalTok{    registry.unclaim("did:lux:alice");  // Should fail but doesn\textquotesingle{}t}
\NormalTok{\}}
\end{Highlighting}
\end{Shaded}

\begin{center}\rule{0.5\linewidth}{0.5pt}\end{center}

\textbf{End of Audit Report}

\end{document}
